(1) 各段线性横振方程为 $\mu_i y_{tt}=Ty_{xx}$，所以 $c_i=\sqrt{T/\mu_i}$、$k_i=\omega/c_i$。固定张力下，行波横向速度与横向力的比定义机械阻抗 $Z_i=T/c_i=\sqrt{T\mu_i}$；阻抗与波数同随线密度平方根增大。
(2) 取左侧 $y_1=Ae^{i(k_1x-\omega t)}+rAe^{i(-k_1x-\omega t)}$，右侧 $y_2=tAe^{i(k_2x-\omega t)}$。在 $x=0$ 处位移连续给 $1+r=t$，横向张力分量连续给 $Tk_1(1-r)=Tk_2t$。解得 $r=(Z_1-Z_2)/(Z_1+Z_2)$、$t=2Z_1/(Z_1+Z_2)$。
(3) 行波的时间平均能流为 $\langle P\rangle=\frac12Z_i\omega^2|A_i|^2$，故 $R_E=|r|^2$、$T_E=(Z_2/Z_1)|t|^2=4Z_1Z_2/(Z_1+Z_2)^2$。直接展开分子，$(Z_1-Z_2)^2+4Z_1Z_2=(Z_1+Z_2)^2$，所以 $R_E+T_E=1$。
(4) $Z_2=Z_1$ 时 $r=0,t=1$；$Z_2\to\infty$ 时 $r\to-1,t\to0$，连接端近于固支；$Z_2\to0^+$ 时 $r\to1,t\to2$，连接端近于自由端。最后一个极限的传输位移虽可达入射的两倍，但其携带能流趋零，绝非能量增益。
